\begin{lemma}
Let $\Omega\subseteq\mathbb{C}$ be open, let $z_0\in\Omega$, let
$g\colon\mathbb{C}\to\mathbb{C}$ and let $L\colon\mathbb{C}\to\mathbb{C}$ be a
continuous $\mathbb{R}$-linear map.  Assume
\begin{enumerate}
\item[(H1)] $g$ is continuous on $\Omega$;
\item[(H2)] $g$ is (real) differentiable at every point of $\Omega-\{z_0\}$;
\item[(H3)] $\mathrm{d}g(z)\to L$ as $z\to z_0$ within $\Omega-\{z_0\}$.
\end{enumerate}
Then $g$ is differentiable at $z_0$ with $\mathrm{d}g(z_0)=L$, and
$z\mapsto\mathrm{d}g(z)$ is continuous at $z_0$ within $\Omega$.

This is the step Section 6 of the paper passes over.  It is what upgrades $C^1$
to $C^k$ across the singular point: by (52), for $|\nu|=m+1\le k$ the exponent
$d-m=k-1+\alpha-m$ is at least $\alpha>0$, so the derivative in question tends
to $0$ at the singular point, and this lemma turns that into genuine
differentiability there.  It is also the continuity at the singular point that
the hypothesis ``$F$ continuous on $U$'' of Lemma 2 of the paper
(paper.tex:538--539) silently presumes.
\end{lemma}
\begin{proof}
Write $h(z)=g(z)-g(z_0)-L(z-z_0)$ for $z\in\Omega$.  By (H1), $h$ is continuous
on $\Omega$ with $h(z_0)=0$; by (H2) and linearity of $L$, $h$ is
differentiable at every $z\in\Omega-\{z_0\}$, with
\begin{equation}\tag{1}
  \mathrm{d}h(z)=\mathrm{d}g(z)-L\qquad(z\in\Omega-\{z_0\}).
\end{equation}
Since $\Omega$ is open and $z_0\in\Omega$, fix $r>0$ with
$B=\{z:|z-z_0|<r\}\subseteq\Omega$.

\emph{Step 1: the radial mean value estimate.}  We claim: if $0<\varepsilon$,
$0<\delta\le r$, and
\begin{equation}\tag{2}
  \bigl\|\mathrm{d}g(w)-L\bigr\|\le\varepsilon\qquad
  \text{for every }w\text{ with }0<|w-z_0|<\delta ,
\end{equation}
then
\begin{equation}\tag{3}
  \bigl\|h(z)\bigr\|\le\varepsilon\,|z-z_0|\qquad\text{for }|z-z_0|<\delta .
\end{equation}

Let $z$ with $|z-z_0|<\delta$.  If $z=z_0$ then $h(z)=0$ and (3) is clear, so
assume $z\neq z_0$.  For $t\in(0,1]$ put $z_t=z_0+t(z-z_0)$, and let
$S_t=[z_t,z]$ be the closed segment joining $z_t$ to $z$.  Every point of
$S_t$ has the form $z_0+s(z-z_0)$ with $s\in[t,1]$, so its distance to $z_0$ is
$s|z-z_0|\in(0,\delta)$; hence $S_t\subseteq B-\{z_0\}\subseteq\Omega-\{z_0\}$
and, by (2) and (1), $\|\mathrm{d}h\|\le\varepsilon$ at every point of $S_t$.
The segment $S_t$ is convex and $h$ is differentiable at each of its points, so
the mean value inequality on a convex set --- if $\|\mathrm{d}h\|\le C$ on a
convex set $S$ then $\|h(y)-h(x)\|\le C|y-x|$ for $x,y\in S$ --- gives
\[
  \bigl\|h(z)-h(z_t)\bigr\|\le\varepsilon\,|z-z_t|
  =\varepsilon(1-t)|z-z_0|\le\varepsilon\,|z-z_0| .
\]
Now let $t\to0^{+}$.  Then $z_t\to z_0$, and $h$ is continuous at $z_0$ with
$h(z_0)=0$, so $h(z_t)\to0$; the left-hand side therefore tends to
$\|h(z)\|$, and (3) follows.

\emph{Step 2: $g$ is differentiable at $z_0$ with derivative $L$.}  Let
$\varepsilon>0$.  By (H3) there is $\delta_0>0$ such that (2) holds with
$\delta=\min\{\delta_0,r\}$: indeed (H3) says precisely that for every
$\varepsilon>0$ the set of $w\in\Omega-\{z_0\}$ with
$\|\mathrm{d}g(w)-L\|\le\varepsilon$ contains the intersection of
$\Omega-\{z_0\}$ with a punctured ball about $z_0$, and shrinking that ball
into $B\subseteq\Omega$ makes every $w$ with $0<|w-z_0|<\delta$ lie in
$\Omega-\{z_0\}$.  By Step 1,
\[
  \bigl\|g(z)-g(z_0)-L(z-z_0)\bigr\|\le\varepsilon\,|z-z_0|
  \qquad\text{for }|z-z_0|<\delta .
\]
As $\varepsilon>0$ was arbitrary, $g(z)-g(z_0)-L(z-z_0)=o(|z-z_0|)$ as
$z\to z_0$, which is exactly the statement that $g$ is differentiable at $z_0$
with $\mathrm{d}g(z_0)=L$.

\emph{Step 3: continuity of the derivative at $z_0$ within $\Omega$.}  By
Step 2, $\mathrm{d}g(z_0)=L$.  Let $\varepsilon>0$ and take $\delta$ as in
Step 2.  For $z\in\Omega$ with $|z-z_0|<\delta$ there are two cases: if
$z\neq z_0$ then $\|\mathrm{d}g(z)-L\|\le\varepsilon$ by (2); and if $z=z_0$
then $\mathrm{d}g(z)-L=L-L=0$.  In both cases
$\|\mathrm{d}g(z)-\mathrm{d}g(z_0)\|\le\varepsilon$.  Hence
$\mathrm{d}g(z)\to\mathrm{d}g(z_0)$ as $z\to z_0$ within $\Omega$, i.e.
$z\mapsto\mathrm{d}g(z)$ is continuous at $z_0$ within $\Omega$.
\end{proof}
